\chapter{The Slip Becomes A Number}

\section{Where the loop fails to close}

The last chapter reframed measurement as closing a loop and reading what came back changed. This chapter follows the
change to its first quantitative form. The place where the loop fails to close has a name in this book --- the
\emph{slip} --- and the slip is where a residue stops being a held remainder and becomes a live reading. The failure is
not a fault to be repaired. The failure is the signal. This section is about learning to read a non-closure as the very
thing the instrument was built to detect.

Set the test up exactly. The machine sends its state around a pass and compares the return to the departure. If the two
agree to the machine's resolution, the loop closed and there is nothing to read. The slip is the case where they do
\emph{not} agree --- where the forward work and the backward work, composed, fail to return the machine to where it
began. Formally it is a residual: the composition of a pass with its own undoing, which ought to be the identity,
instead leaves a remainder. That remainder is the slip. It is the exact amount by which the machine's own round-trip
fails to be faithful, and it is a quantity the machine can compute directly, by running the loop and measuring the gap.

There is a change of character here worth naming, because it is the hinge of the whole chapter. Until now the residue
was a \emph{static} thing --- a tail held in the representation, a leftover carried across the seam, a remainder at
rest. At the slip the same residue turns \emph{active}: it stops being a value the machine stores and becomes a value
the machine \emph{reads off the loop's failure to close}. The distinction is the distinction between a quantity you hold
and a quantity you measure. The slip is the residue caught in the act --- not sitting in a cell but showing itself as
the discrepancy between where the loop should have returned and where it did. The machine has, in effect, discovered
that its own imperfect round-trip is an instrument: the gap it leaves is a reading of something real.

And the something real is exactly what the whole book is chasing. A loop that closed perfectly would be a machine with
no residue, and a machine with no residue would have nothing to measure --- it would be a perfect, empty mirror. The
slip is the loop confessing that it carries a remainder it cannot dispose of, and that confession is the measurement.
The failure to close is not the instrument breaking; it is the instrument working. This is the deepest inversion in the
early book: the machine measures by the imperfection of its own loop, and a machine whose loops all closed clean would
be an instrument that could read nothing at all.

In the two verbs the slip is a tange the loop performs on itself. The round-trip tanges the returned state against the
departed --- decides whether they are the same --- and where the decision comes back \emph{not-same}, the difference it
isolates is the slip. What comes back the same the machine funges away, unread; the slip is the residue the funge cannot
absorb, surfaced as a live discrepancy. So the slip is the machine's two verbs turned on the loop's own faithfulness: it
tanges the failure apart from the success and reads the failure as the signal. What that signal is worth --- how the
slip becomes not just a discrepancy but a bracketed number --- is the next section's.

\section{Bracketed, not guessed}

The slip is a discrepancy the machine can compute. This section makes it a \emph{number}, and the whole of the section
is about the manner in which it does so, because the manner is the book's entire honesty in miniature. The machine does
not fit the slip, tune it, or guess it. It \emph{brackets} it --- bounds it from below and from above by direct
computation --- and reports the bracket. A slip guessed would be worthless; a slip bracketed is a reading. The
difference between the two is the difference between this instrument and a crank's.

Watch the bracketing. The machine does not ask ``what is the slip?'' and produce a point. It asks two decidable
questions --- ``is the slip at least this much?'' and ``is the slip at most that much?'' --- and answers each by
computation. A lower wall it can certify: the slip is provably no smaller. An upper wall it can certify: the slip is
provably no larger. Between the two walls sits the value, bounded on both sides, with the machine under no obligation to
say where between them it lies. This is the bisection method, the oldest honest way to corner a value a machine cannot
evaluate directly~\cite{trefethen1997}: bracket it, halve the gap, and halve it again, each halving a decidable
comparison that tightens one wall or the other. The slip becomes a number by being cornered, not by being named.

The point that a bracket is \emph{earned} where a guess is merely asserted deserves its full weight, because it is why
the machine can be trusted at the end. A fitted value --- a number chosen because it makes something come out right ---
carries no certificate; you cannot ask it what it rests on. A bracketed value carries two certificates, one per wall,
each a computation the machine actually ran and can rerun. The machine never says ``the slip is this''; it says ``the
slip is provably above here and provably below there, and here is the computation for each.'' Guessing would let the
machine pretend to a precision it has not earned; bracketing binds it to exactly the precision its computations
certify, and not a hair more. The bracket is the honest form of a number the machine cannot compute exactly, and the
slip is the first place the machine is forced to use it.

There is a floor to the bracketing, and it must be named now because it governs the rest of the book, though the next
chapter is where it is built. The halving cannot continue forever. At some depth the two walls fall closer together than
the machine can resolve --- the comparison that would tighten them further returns \emph{indistinguishable}, because the
gap has dropped below the machine's own resolution. There the bisection must stop. It does not stop because the value is
found; it stops because the machine has run out of decidable distinctions between the walls. The bracket at that depth
is the honest reading: as tight as the machine's resolution allows, and no tighter. To halve past it would be to invent
distinctions the machine cannot decide --- to guess, dressed as computation --- which is exactly the dishonesty the
whole method exists to refuse.

In the two verbs the bracketed slip is a funge held to a finite width. The machine tanges the slip apart from below and
from above --- decides it larger than the lower wall, smaller than the upper --- and funges together, into a single
interval, every value the two decided walls cannot separate. A guess would funge the interval shut on a point the
machine never decided; the bracket keeps the funge open, exactly as wide as the decided walls require. The slip is a
number now, but a number of the machine's honest kind: an interval it earned, not a point it wished for. What it means
for a symbol to have become this kind of number --- an instrument reading rather than a written mark --- is the last
section's.

\section{The reading}

The machine has a bracketed slip: a value cornered between two computed walls. This closing section names the
transformation that has quietly happened over the chapter, because it is a change in the meaning of the word
\emph{number} itself. Up to the slip, a number was a symbol --- a codeword the machine wrote, a finite condition it
specified, a mark in its representation. At the slip, a number becomes a \emph{reading} --- a value the machine took off
its own working, cornered by computation, carrying its own certificate. The chapter's title is literal: the slip
becomes a number, and in becoming one it turns number from something written into something read.

Feel the difference between the two. A written number is authored: the machine puts it there, and it means whatever the
machine defined it to mean. A read number is \emph{found}: the machine did not author it but cornered it, and it means
the thing the machine's loop failed to close by. The written number of the early chapters was a representation --- a
value the machine held. The read number of this chapter is a measurement --- a value the machine took. This is the
first genuine instrument reading in the book: not a symbol the machine wrote down but a bracket it computed off the
imperfection of its own round-trip. The machine has stopped merely representing values and started measuring one.

What makes a reading a reading, and not just another written number, is that it answers to something outside the
machine's choice. The machine chose its codewords; it did not choose its slip. The slip is whatever the loop's failure
to close makes it, and the machine can only corner it, not author it. This is the sense in which the reading is
\emph{objective} within the machine's finite world: it is the one number the machine did not get to decide, only to
bracket. Every later reading the instrument takes --- the count of turns, the cost of a proof, the coupling at the very
end --- is this same kind of thing: a value the machine finds by cornering, not a value it writes by choosing. The slip
is the training case for reading itself, the first number the machine measured rather than made.

And the manner of the reading fixes the manner of every reading to come. Because the slip came as a bracket, every
reading after it comes as a bracket: the machine, having learned that its honest numbers are intervals cornered between
computed walls, does not later pretend to points. The coupling at the end will be a bracket for exactly the reason the
slip is one --- because it is a value the machine cornered and could not author, floored at the machine's own
resolution. The chapter that taught the machine to read taught it, in the same stroke, that a reading is an interval.
The instrument reads in brackets from here to the end, and it learned the habit here, on the first value it ever
measured rather than made.

In the two verbs the reading is the tange become objective. Every tange before this the machine performed by its own
choice of characteristic; the slip is the tange the loop performs on the machine, a difference the machine did not
choose but must read. The machine funges its written numbers as it likes --- pools and names them by codewords of its
own devising --- but the read number it can only corner, funging the indistinguishable values into a bracket it did not
author. So the slip is the first place the machine's two verbs answer to something beyond the machine: a difference it
did not choose, cornered into an interval it did not pick, and read as the instrument's first true number. The next
chapter takes this bracket seriously as an object in its own right --- the bound, not the point, as the thing the
machine owns.
